\begin{enumerate}
\item[\textbf{(b)}] Using (a), find a UMVUE for $\alpha$.
\end{enumerate}

\textbf{5.3.18.} Let $X_1, \ldots, X_n$ be a random sample from a population with density function:
$$f(x) = \begin{cases} \dfrac{1}{\theta}, & -\dfrac{\theta}{2} \leq x \leq \dfrac{\theta}{2}, \quad \theta > 0 \\ 0, & \text{otherwise.} \end{cases}$$

Show that $\left(\min\limits_{1 \leq i \leq n} X_i, \max\limits_{1 \leq i \leq n} X_i\right)$ is sufficient for $\theta$.

\textbf{5.3.19.} Let $X_1, \ldots, X_n$ be a random sample from a $G(1, \beta)$ distribution.
\begin{enumerate}
\item[\textbf{(a)}] Show that $U = \sum_{i=1}^{n} X_i$ is a sufficient statistic for $\beta$.
\item[\textbf{(b)}] The following is a random sample from a $G(1, \beta)$ distribution:
\end{enumerate}

\begin{center}
\begin{tabular}{cccccccccc}
0.3 & 3.4 & 0.4 & 1.8 & 0.7 & 1.0 & 0.1 & 2.3 & 3.7 & 2.0 \\
0.3 & 3.7 & 0.1 & 1.3 & 1.2 & 3.3 & 0.2 & 0.2 & 1.3 & 0.6 \\
\end{tabular}
\end{center}

Find a sufficient statistic for $\beta$.

\textbf{5.3.20.} Show that $X_1$ is not sufficient for $\mu$, if $X_1, \ldots, X_n$ are a sample from $N(\mu, 1)$.

\textbf{5.3.21.} Let $X_1, \ldots, X_n$ be a random sample from the truncated exponential distribution with pdf:
$$f(x) = \begin{cases} e^{\theta - x}, & x > \theta \\ 0, & \text{otherwise.} \end{cases}$$

Show that $X_{(1)} = \min(X_i)$ is sufficient for $\theta$.

\textbf{5.3.22.} Let $X_1, \ldots, X_n$ be a random sample from a distribution with pdf:
$$f(x) = \begin{cases} \theta x^{\theta-1}, & 0 < x < 1, \quad \theta > 0 \\ 0, & \text{otherwise.} \end{cases}$$

Show that $U = Y_1, \ldots, Y_n$ is a sufficient statistic for $\theta$.

\textbf{5.3.23.} Let $X_1, \ldots, X_n$ be a random sample from a Rayleigh distribution with pdf:
$$f(x) = \begin{cases} \dfrac{2x}{\alpha} e^{-x^2/\alpha}, & x > 0 \\ 0, & \text{otherwise.} \end{cases}$$

Show that $\sum_{i=1}^{n} X_i^2$ is sufficient for the parameter $\alpha$.

\section{A method of finding the confidence interval: pivotal method}

In the previous sections, we studied methods for finding point estimators for the population parameters. In general, the estimates will differ from the true parameter values by varying amounts depending on the sample values obtained. The point estimates do not convey any measure of reliability.

Now, we discuss another type of estimation, called an \textit{interval estimation}. Although point estimators are useful, interval estimators convey more information about the data that are used to obtain the point estimate. The purpose of using an interval estimator is to have some degree of confidence of securing the true parameter. For an interval estimator of a single parameter $\theta$, we will use the random sample to find two quantities $L$ and $U$ such that $L < \theta < U$ with some probability. Because $L$ and $U$ depend on the sample values, they will be random. This interval $(L, U)$ should have two properties: (1) $P(L < \theta < U)$ is high, that is, the true parameter $\theta$ is in $(L, U)$ with high probability, and (2) the length of the interval $(L, U)$ should be relatively narrow on average.

In summary, interval estimation goes a step beyond point estimation by providing, in addition to the estimating interval $(L, U)$, a measure of one's confidence in the accuracy of the estimate. Interval estimators are called \textit{confidence intervals} and the limits are called $U$ and $L$, the \textit{upper} and \textit{lower confidence limits}, respectively. The associated levels of confidence are
